\section{Results}

All five sub-hypotheses passed validation gates with metrics exceeding targets, demonstrating mechanistic feasibility of the three-component system at proof-of-concept scale.

\subsection{Instrumentation Feasibility (h-e1)}

\textbf{Overhead:} 0.21\% (target: <10\%)
\textbf{Capture Rate:} 100\% (target: $\geq$95\%)
\textbf{Events Captured:} 100 over simulated 6-month window (target: $\geq$100)

Load-time instrumentation added negligible latency (<2ms per dataset load). SHA256 hashing amortized via lazy caching; async queue writes minimized blocking I/O. Overhead significantly below 10\% tolerance suggests production deployment viable without performance degradation.

\subsection{Health Metrics Precision (h-m1)}

\textbf{Precision:} 88.3\% (target: $\geq$60\%)
\textbf{Recall:} 100\% (target: $\geq$80\%)
\textbf{Weighted Score Threshold:} $\geq$2.5

Automated health metrics exceeded precision target by 47\%. Weighted scoring (velocity=2.0, emergence=1.5, issue\_ratio=1.0) achieved strong discrimination on 1000-dataset simulation. Zero false negatives (100\% recall) validates that all true deprecation candidates surfaced by metrics.

Confusion matrix analysis shows declining usage velocity (velocity < 0.3) was strongest predictor; issue ratio contributed less signal but improved precision for edge cases (maintained datasets with declining usage due to successor availability rather than quality issues).

\subsection{Context Inference Accuracy (h-m2)}

\textbf{Accuracy:} 100\% (target: $\geq$70\%)
\textbf{User Override Rate:} 27.95\% (target: <50\%)

Pattern-based context inference achieved perfect accuracy on synthetic validation dataset. 72\% acceptance rate (100\% - 27.95\% override) suggests context-aware recommendations add value over linear version succession.

Accuracy expected to regress to 70-90\% on real-world usage patterns (synthetic patterns cleaner than production data). Override rate provides safety valve: users can correct misclassified context, preventing incorrect successor recommendations.

\subsection{Migration Planning Completeness (h-m3)}

\textbf{Impact Completeness:} 100\% (target: $\geq$95\%)
\textbf{Schema Accuracy:} 100\% (target: $\geq$80\%)

Transitive closure analysis on 4-node dependency graph identified all affected entities (datasets, downstream models, training pipelines). Field-level schema diff detected all breaking changes (column renames, type changes, missing fields) automatically.

Validated on small-graph scale (4 nodes); 50k-node scalability inferred but not tested. NetworkX graph algorithms scale to thousands of nodes with sub-second query times; production validation required for 60k-dataset HuggingFace scale.

\subsection{Full-Stack Telemetry (h-m4)}

\textbf{Capture Rate:} 100\% (95\% CI: 99.05\%-100\%) (target: $\geq$95\%)
\textbf{Overhead:} 5\% (target: <10\%)
\textbf{Event Completeness:} 100\% (target: $\geq$95\%)

Integrated system (health metrics + context graphs + instrumentation) maintained high capture rate with overhead below tolerance. Async queue design scaled to 500 simulated user sessions without blocking. Wilson confidence interval lower bound (99.05\%) exceeds 95\% target, validating statistical reliability.

Overhead higher than h-e1 baseline (5\% vs 0.21\%) due to full telemetry stack (metadata capture, dependency graph queries, health metric computation), but still well below 10\% tolerance.

\subsection{Cross-Hypothesis Patterns}

All five sub-hypotheses (h-e1, h-m1, h-m2, h-m3, h-m4) passed gates, validating mechanistic integration:

\begin{itemize}
\item \textbf{Component 1 (Health Metrics):} Automated detection feasible with 88.3\% precision
\item \textbf{Component 2 (Context Graphs):} Task-aware recommendations feasible with 100\% accuracy on synthetic patterns
\item \textbf{Component 3 (Instrumentation):} Adoption tracking feasible with <5\% overhead and 100\% capture
\end{itemize}

Results validate infrastructure layer (measurement + mechanisms work) but NOT efficacy layer (adoption lift untested---Phase 5 baseline comparison was skipped).

\subsection{Surprising Findings}

Instrumentation overhead significantly lower than expected (0.21\% baseline, 5\% full stack). SHA256 hashing expected to dominate overhead; lazy caching and background queue writes effectively amortized cost. Perfect scores (100\% recall, 100\% context accuracy) likely optimistic due to synthetic data; real-world accuracy expected to regress to 70-95\% based on pattern cleanliness assumptions.
